\documentclass[11pt,a4paper]{article}
\usepackage{turkos}

\begin{document}
\phaseheader{15}{Hardening, Polish and Release: Turk-OS 1.0}{Make it survive hostile input, measure it, document it, ship it, and choose where to go next}{96}{100}

\ataglance{3--4 weeks for the core, then open-ended}{4}{Phases 0--14 complete: Turk-OS boots to a shell with pipes, files and processes}{A
system call fuzzer that runs for an hour without a kernel panic, an audited user/kernel boundary, kernel stack
protection and read-only kernel code, users and file permissions, a benchmark suite with one measured optimisation,
an architecture document, a reproducible \code{make release} that produces a bootable ISO, the \code{v1.0} tag, and a
written reflection on your design. After that, a choice of tracks beyond 100\,\%.}

\section{Why this phase}

A system isn't finished when its features work on the inputs you tried. It is finished when it survives inputs
nobody intended, when you know how fast it is and why, when someone else can build and understand it, and when you
can explain every design choice in it. That is the difference between a demo and version 1.0, and this phase is about
closing that gap.

It is also a good moment to step back. \MOS chapter 12, written by an author who designed several operating systems,
is about exactly the decisions you have been making for nine months: interfaces, structure, mechanism versus policy,
performance, and managing a large software project. Reading it now, with your own kernel in mind, is very different
from reading it at the start. The phase ends with a written comparison of Turk-OS against MINIX 3 and Linux, and with
a list of optional tracks for going further, one of which connects Turk-OS to the RISC-V core you are building in
hardware.

\section{Concepts you need}

\subsection{The kernel is the trusted computing base}

\MOS\,\S9.2.2 defines the \textbf{trusted computing base} (TCB): the hardware and software that must be correct for
the system's security to hold. In Turk-OS that is the entire kernel. Every system call is an entry point into the TCB
controlled by a possibly hostile program, which is why Phase 10's rule (never trust a user pointer) was only the
start. \MOS\,\S9.7 catalogues how kernels get broken, and several items apply directly to code you have written:
\textbf{buffer overflows} when a length from the user isn't checked; \textbf{integer overflows} when
\code{offset + length} wraps around; \textbf{null-pointer dereferences} in the kernel; and \textbf{time-of-check to
time-of-use} (TOCTOU) races, where the kernel checks a value in user memory and then reads it again after the user has
changed it. The cure for the last one is to copy first and check the copy.

\subsection{Protection: who may do what}

Protection needs a model of \emph{subjects} (users, processes), \emph{objects} (files, devices) and \emph{rights}
(read, write, execute). \MOS\,\S9.3 and \OSC chapter 17 describe the two ways to store it: \textbf{access control
lists} attached to objects, and \textbf{capabilities} held by subjects. UNIX uses a compact ACL, nine permission bits
for owner, group and others in each inode, plus the rule that user ID 0 (root) may do anything. TurkFS already stores
the mode, user ID and group ID in every inode (Phase 13), so what is left is to give processes an identity and check it.

\subsection{Measuring before optimising}

\MOS\,\S12.4 gives the rules for performance work: find out what is actually slow, optimise the common case, use
caching and hints, and exploit locality. None of that works without measurements. The x86 \code{rdtsc} instruction
reads a cycle counter, which is precise enough to time a single system call. The classic operating-system
microbenchmarks (in the style of \emph{lmbench}) measure the null system call, a context switch between two processes
passing a token through pipes, process creation, and file-system throughput.

\subsection{Design, looking back}

\MOS\,\S12.2--12.3 lists principles you can now test against your own code: \textbf{simplicity} and
\textbf{completeness} of the interface, \textbf{orthogonality} (independent concepts that combine freely, like
\code{fork} and \code{exec}), \textbf{separation of mechanism and policy} (your scheduler), static versus dynamic
structures (MINIX's fixed tables versus your heap), and top-down versus bottom-up implementation (you went bottom-up).
\S12.5 adds Brooks's observations about software projects, including the \emph{second-system effect}: the temptation to
put everything you left out of version 1 into version 2.

\section{Reading plan}

\begin{readingplan}
\must & \MOS & Ch.~12 \emph{Operating System Design}: nature of the problem, interface design, implementation, performance, project management, trends & 981--1030 & Read all of it, with Turk-OS in mind. \\
\must & \MOS & \S9.1--9.3 \emph{The Security Environment}, \emph{Operating Systems Security}, \emph{Controlling Access to Resources} & 595--611 & Threats, the TCB, protection domains, ACLs and capabilities. \\
\must & \MOS & \S9.7 \emph{Exploiting Software} (buffer overflows, format strings, dangling and null pointers, integer overflows, TOCTOU) & 639--657 & The bugs your audit is looking for. \\
\should & \MOS & \S9.10 \emph{Defenses} (especially code signing and jailing) & 684--703 & What operating systems add on top of correct code. \\
\should & \OSC & \S16.1--16.2 \emph{The Security Problem}, \emph{Program Threats}; \S16.5--16.6 \emph{User Authentication}, \emph{Implementing Security Defenses} & 621--634, 648--662 & A second view, with passwords and defences in depth. \\
\should & \OSC & Ch.~17 \emph{Protection}, \S17.1--17.6 (principles, rings, domains, access matrix) & 667--682 & The formal model behind UNIX permissions. \\
\should & \OSDI & \S5.4--5.5 \emph{Security}, \emph{Protection Mechanisms} & 526--548 & Compact coverage of authentication, ACLs, capabilities and covert channels. \\
\should & \OSDI & \S1.5 \emph{Operating System Structure}; \S2.5.1 \emph{The Internal Structure of MINIX 3} & 42--51, 112--116 & Monolithic versus microkernel, for your design reflection. \\
\should & \MOS & \S10.1--10.8 \emph{Case Study: UNIX, Linux, and Android} (skim) & 713--848 & Compare each subsystem with yours. \\
\could & \OSC & Ch.~20 \emph{The Linux System} & 775--819 & Another view of Linux for the comparison. \\
\end{readingplan}

\begin{minixbox}
MINIX 3's central design goal is reliability, and \OSDI\,\S2.5.1 explains how its structure serves it. Drivers run
as separate user processes with only the privileges they need, so a buggy driver can't overwrite the kernel. A
special process, the \textbf{reincarnation server}, watches the drivers and servers and restarts any that crash
(\OSDI pp.~114 and 599). In Turk-OS, every driver bug is a kernel bug. Keep this contrast in mind for Task~15.8, and
look at Track F below if you want to try the MINIX approach on one of your own drivers.
\end{minixbox}

\section{Build it}

\task{Fuzz the system call interface}
Write \code{user/bin/fuzz}: a program that makes random system calls forever, with a seed you can print and replay.
Mix the arguments: small integers, huge integers, negative numbers, pointers to valid buffers, pointers into the kernel
(\code{0xC0000000} and up), unmapped addresses, addresses one byte before the end of a page, descriptors that exist and
ones that don't, and random strings as paths. Run several copies at once together with a stress script that forks,
execs, pipes and writes files in parallel. Every kernel panic is a bug to fix; every process killed by a signal is
fine.
\verify{Four fuzzers and the stress script run for one hour with zero kernel panics and no growth in memory use
after they are stopped.}

\task{Audit the boundary}
Go through every system call and write down, in a table in \code{docs/}, each argument and how it is validated. Look
specifically for these bugs, because the fuzzer is unlikely to find all of them.

\begin{center}
\small\renewcommand{\arraystretch}{1.25}
\begin{tabularx}{\linewidth}{@{}l X@{}}
\toprule
\tkth{Bug class} & \tkth{Where to look in Turk-OS} \\
\midrule
missing pointer check & any use of a user pointer outside \code{copy_from_user}/\code{copy_to_user}/\code{strncpy_from_user} \\
integer overflow & \code{offset + len} in \code{read}/\code{write}/\code{lseek}, \code{brk} arguments, ELF segment sizes, \code{argv} totals \\
TOCTOU & paths or structures read twice from user memory; copy once into the kernel, then validate the copy \\
unchecked index & descriptor numbers, signal numbers, system call numbers, directory entry indices \\
information leak & kernel stack or heap bytes copied to the user in padding of \code{struct stat} or \code{dirent}; zero structures before filling them \\
resource exhaustion & unlimited processes, descriptors, pipe buffers or path lengths; every resource needs a limit \\
\bottomrule
\end{tabularx}
\end{center}

\task{Harden the kernel itself}
Compile the kernel with \code{-fstack-protector-strong}. GCC then inserts a canary check in functions with local
arrays and needs two symbols from you: \code{__stack_chk_guard}, set to an unpredictable value at boot (for example
from \code{rdtsc}), and \code{__stack_chk_fail}, which panics. Map the kernel image with 4\,KiB pages so \code{.text} and
\code{.rodata} are read-only (with CR0.WP set, a stray write faults at once). Make sure the page at address 0 is never
mapped in any process. Finally, make crash reports readable: generate a symbol table from \code{i686-elf-nm} at build
time, load it as a boot module, and print function names in the backtrace of every panic.
\verify{A deliberate overflow of a local buffer in a test system call panics with ``stack smashing detected'' and a
symbolised backtrace; a write to a function's code faults.}

\task{Users and permissions}
Give each task a user ID and a group ID, inherited through \code{fork}. Add \code{getuid} (24), \code{setuid} (23),
\code{getgid} and \code{setgid}; only root may change identity. Check permissions in the VFS against the inode's mode,
owner and group on \code{open}, \code{execve}, \code{create}, \code{mkdir} and \code{unlink}, with root bypassing the
checks. Add \code{/etc/passwd} with lines of the form \code{name:uid:gid:home:shell}, and a \code{login} program that
\code{init} runs on the console, which looks up the user, changes to their identity and home directory, and runs their
shell. If you add passwords, store only salted hashes, never the passwords themselves.
\verify{An ordinary user can't read a file with mode \code{0600} owned by root, can't write to \code{/bin}, and can't
kill root's processes; root can do all three.}

\task{Benchmarks, and one real optimisation}
Write \code{user/bin/bench}, which reports the median of many runs for each measurement below, and put the results in
a table in your documentation. Then pick the worst number, find the cause with measurements, fix it, and record the
before and after values.

\begin{center}
\small\renewcommand{\arraystretch}{1.25}
\begin{tabularx}{\linewidth}{@{}l X@{}}
\toprule
\tkth{Benchmark} & \tkth{How} \\
\midrule
null system call & time 100\,000 calls to \code{getpid} with \code{rdtsc} \\
context switch & two processes pass a byte back and forth through two pipes; divide by the number of switches \\
process creation & \code{fork} + \code{exit} + \code{wait}; and \code{fork} + \code{exec} + \code{wait} of a tiny program \\
file throughput & write and read a 16\,MiB file on TurkFS, with a cold and a warm buffer cache \\
boot time & \code{rdtsc} in the loader and at the first shell prompt, printed on the serial port \\
\bottomrule
\end{tabularx}
\end{center}

Typical optimisations at this point are global pages for kernel mappings (no TLB flush of kernel entries on a CR3
switch), multi-sector disk transfers and read-ahead, a hash table in the buffer cache, and avoiding the double copy
through a kernel buffer in pipes. Run the same benchmarks in a Linux system for perspective; the point is to
understand the differences, not to win.

\task{Write the documentation}
Write an architecture document (a LaTeX file next to this plan, or Markdown): an overview diagram like the one in the
roadmap; each subsystem in a page or two; the virtual and physical memory maps; the system call table with numbers,
arguments and errors; the lock order; the TurkFS on-disk format; how to build, run, test and debug; and known
limitations. Add a \code{README} that gets a newcomer from a clean clone to a running shell, a \code{CHANGELOG}, and a
license (MIT matches your other open-source work).

\task{Release v1.0}
Add \code{make release}: build the kernel, the user programs, the root archive and a TurkFS image populated with
\code{/bin}, \code{/etc} and a home directory, and a GRUB ISO whose \code{grub.cfg} loads the kernel and its modules.
Check that the release builds from a clean clone (\code{git clone} into a temporary directory, then \code{make release}),
boot the ISO in QEMU, and tag the commit \code{v1.0}. If you want to try real hardware, an older PC with legacy BIOS (or
UEFI with CSM) can boot the ISO from a USB stick.

\begin{warnbox}[Writing a USB stick]
\code{sudo dd if=build/turkos.iso of=/dev/sdX bs=4M} erases the target device completely. Run \code{lsblk} before and
after plugging the stick in, make sure \code{sdX} is the stick and not your laptop's disk, and never guess.
\end{warnbox}
\verify{A fresh clone builds with one command and boots from ISO to the login prompt. \code{make test} passes on the
tagged commit.}

\task{Defend your design}
Write a reflection of three to five pages, as if for a thesis committee. For each subsystem (boot, memory, processes,
scheduling, synchronisation, system calls, file system, drivers), state what Turk-OS does, how MINIX 3 does it
(\OSDI) and how Linux does it (\MOS chapter 10), and what you would change with hindsight. Use \MOS chapter 12's
principles as the yardstick, and end with the one decision you are proudest of and the one you would reverse first.
Presenting it to your advisor or classmates is the best test of whether you really understand your own system.

\section{Beyond 100\,\%: choose a track}

Turk-OS 1.0 is a complete small system. These tracks extend it in different directions; each is a project of several
weeks to several months, and they can be combined. Pick by interest, not by size.

\begin{center}
\small\renewcommand{\arraystretch}{1.35}
\begin{longtable}{@{}>{\raggedright\arraybackslash}p{0.19\linewidth} >{\raggedright\arraybackslash}p{0.46\linewidth} >{\raggedright\arraybackslash}p{0.29\linewidth}@{}}
\toprule
\tkth{Track} & \tkth{What you build} & \tkth{Where to read} \\
\midrule
\endfirsthead
\toprule
\tkth{Track} & \tkth{What you build} & \tkth{Where to read} \\
\midrule
\endhead
\bottomrule
\endlastfoot
A.\ Graphics and Turkish text & Request a linear framebuffer through the Multiboot video fields, draw text with a PSF font that covers the whole Turkish alphabet (ğ, ş, ı, İ, Ğ, Ş), add a PS/2 mouse (IRQ 12), then a simple window system & \MOS\,\S5.6 (user interfaces); OSDev Wiki (VESA, PSF, mouse) \\
B.\ Multiprocessor (SMP) & Parse the ACPI MADT, program the local APIC and I/O APIC, start the other CPUs (INIT--SIPI--SIPI), per-CPU data and run queues, TLB shootdown; your Phase 9 spinlocks finally spin & \MOS\,\S8.1 (520--544); \OSTEP Ch.~10 (97--108); \OSC\,\S5.5 (220--227) \\
C.\ Networking & PCI enumeration, an e1000 or RTL8139 driver in QEMU, Ethernet, ARP, IPv4, ICMP (\code{ping}), UDP, then a minimal TCP and a socket API & \OSC Ch.~19 (733--768); \MOS\,\S8.3 (566--587); OSDev Wiki and RFCs \\
D.\ 64-bit port & Long mode, four-level paging, Multiboot2, \code{syscall}/\code{sysret}, the x86-64 calling convention & \OSC\,\S9.6 (379--383); Intel SDM Vol.~3A \\
E.\ RISC-V port & Separate the architecture layer (boot, traps, context switch, page tables), port it to RISC-V with Sv32 paging (two levels, like x86), boot on QEMU's \code{virt} machine, then on your own Timur-RV32IMC core once it has supervisor mode and a timer & RISC-V Privileged Specification; xv6-riscv as a reference \\
F.\ Microkernel experiment & Move the ATA driver (or TurkFS) into a user-space server that talks to the kernel with synchronous messages, as MINIX does, and measure the cost & \OSDI\,\S2.5, \S3.4 (112--125, 252--261); \MOS\,\S1.7.3 (65--68) \\
G.\ Virtualisation & Understand what QEMU and KVM do for you, run Turk-OS with \code{-enable-kvm}, and compare the benchmarks & \OSC Ch.~18 (701--729); \MOS Ch.~7 (471--516) \\
\end{longtable}
\end{center}

Track E deserves a note. Porting forces you to find every place where x86 leaked into supposedly generic code, which
improves the whole kernel. And running your own operating system on a processor you designed yourself is a rare
thing to be able to say about a student project, and a natural thesis topic.

\section{Milestone}

\begin{milestonebox}[Turk-OS 1.0]
\begin{checklist}
  \item One hour of parallel fuzzing and stress testing without a kernel panic or a memory leak.
  \item Boundary audit table complete; stack protector, read-only kernel code and symbolised backtraces in place.
  \item Users, \code{login} and file permissions enforced.
  \item Benchmark table written, with one optimisation measured before and after.
  \item Architecture document, README, system call reference and CHANGELOG written.
  \item \code{make release} works from a clean clone; the ISO boots; \code{v1.0} is tagged.
  \item The design reflection is written, and ideally presented.
\end{checklist}
\end{milestonebox}

\begin{questionbox}
\begin{enumerate}
  \item Which parts of Turk-OS trust input from user programs, and how is each validated?
  \item Give an example of a TOCTOU bug that could exist in \code{execve}, and explain how copying first prevents it.
  \item Where would MINIX 3 put your ATA driver, and what would change about a driver bug's consequences?
  \item What is the most expensive operation you measured, and where does the time go?
  \item Which of \MOS chapter 12's principles did Turk-OS follow best, and which did it violate most?
  \item What would break first if Turk-OS ran on two CPUs tomorrow?
\end{enumerate}
\end{questionbox}

\section{Closing}

You started this plan knowing basic C and no assembly. Turk-OS 1.0 boots on a PC, protects processes from each other,
schedules them preemptively, manages virtual memory with copy-on-write, stores files in a real on-disk format, and runs
a shell with pipes. More useful than the code is what you now understand: what happens between pressing a key and a
character appearing, between \code{fork} and \code{exit}, between \code{write} and the disk. The five books on your
shelf will read differently from now on, because every chapter describes something you have built or decided not to
build. Keep the journal going, and pick a track.

\booklegend
\end{document}
